% GENERATED FILE. Edit canonical lessons, then run npm run generate:books.
% canonical-source-hash: fnv1a64:e6965ed77e6bff38
% canonical-lessons: SA-C255-siksa, SA-C255-jnanam, SA-C255-abhyasah, SA-C255-grhakaryam, SA-C255-uttaram

\chapter{Education, Knowledge, Practice, Homework, Answer}
\label{ch:sa-education-knowledge-practice-homework-answer}
\hlchaptermodality{255}

\begin{chapteropening}
\textbf{By the end of this chapter:} \emph{I can say education, knowledge, practice, homework and an answer.}

Complete the last lesson of chapter 255: I can say education, knowledge, practice, homework and an answer.
\end{chapteropening}

\section[\texorpdfstring{\sk{शिक्षा} (śikṣā)}{शिक्षा (śikṣā)}]{\sk{शिक्षा} (śikṣā) — education}

\label{lesson:SA-C255-siksa}



\begin{quote}
Before the new one: say the Sanskrit for dinner, then the Sanskrit for cooked food.
\end{quote}

\subsection*{You'll want to know: \sk{शिक्षा}}
\textbf{\sk{शिक्षा}} — \emph{śikṣā} — \textquotedblleft{}education\textquotedblright{}.

\begin{grammarlens}[title={What you've built}]
One more word for this chapter.
\end{grammarlens}

\subsection*{Your turn}
\begin{itemize}
\raggedright
  \item \emph{Say it:} \emph{śikṣā}
  \item \emph{Say it:} \emph{śikṣā}, once more
  \item \emph{Say it:} say \emph{śikṣā}
  \item \emph{Recall:} say the Sanskrit for dinner, then the Sanskrit for cooked food, then say \emph{śikṣā} again
\end{itemize}

\begin{tcolorbox}[breakable,colback=teal!4,colframe=teal!35!black,title={Before you move on}]
What does \sk{शिक्षा} mean? (\textquotedblleft{}Education\textquotedblright{}.) Say it once more.
\end{tcolorbox}
\section[\texorpdfstring{\sk{ज्ञानम्} (jñānam)}{ज्ञानम् (jñānam)}]{\sk{ज्ञानम्} (jñānam) — knowledge}

\label{lesson:SA-C255-jnanam}



\begin{quote}
Before the new one: say the Sanskrit for cooked food, then the Sanskrit for education.
\end{quote}

\subsection*{You'll want to know: \sk{ज्ञानम्}}
\textbf{\sk{ज्ञानम्}} — \emph{jñānam} — \textquotedblleft{}knowledge\textquotedblright{}.

\begin{grammarlens}[title={What you've built}]
One more word for this chapter.
\end{grammarlens}

\subsection*{Your turn}
\begin{itemize}
\raggedright
  \item \emph{Say it:} \emph{jñānam}
  \item \emph{Say it:} \emph{jñānam}, once more
  \item \emph{Say it:} say \emph{jñānam}
  \item \emph{Recall:} say the Sanskrit for cooked food, then the Sanskrit for education, then say \emph{jñānam} again
\end{itemize}

\begin{tcolorbox}[breakable,colback=teal!4,colframe=teal!35!black,title={Before you move on}]
What does \sk{ज्ञानम्} mean? (\textquotedblleft{}Knowledge\textquotedblright{}.) Say it once more.
\end{tcolorbox}
\section[\texorpdfstring{\sk{अभ्यासः} (abhyāsaḥ)}{अभ्यासः (abhyāsaḥ)}]{\sk{अभ्यासः} (abhyāsaḥ) — practice}

\label{lesson:SA-C255-abhyasah}



\begin{quote}
Before the new one: say the Sanskrit for education, then the Sanskrit for knowledge.
\end{quote}

\subsection*{You'll want to know: \sk{अभ्यासः}}
\textbf{\sk{अभ्यासः}} — \emph{abhyāsaḥ} — \textquotedblleft{}practice\textquotedblright{}.

\begin{grammarlens}[title={What you've built}]
One more word for this chapter.
\end{grammarlens}

\subsection*{Your turn}
\begin{itemize}
\raggedright
  \item \emph{Say it:} \emph{abhyāsaḥ}
  \item \emph{Say it:} \emph{abhyāsaḥ}, once more
  \item \emph{Say it:} say \emph{abhyāsaḥ}
  \item \emph{Recall:} say the Sanskrit for education, then the Sanskrit for knowledge, then say \emph{abhyāsaḥ} again
\end{itemize}

\begin{tcolorbox}[breakable,colback=teal!4,colframe=teal!35!black,title={Before you move on}]
What does \sk{अभ्यासः} mean? (\textquotedblleft{}Practice\textquotedblright{}.) Say it once more.
\end{tcolorbox}
\section[\texorpdfstring{\sk{गृहकार्यम्} (gṛhakāryam)}{गृहकार्यम् (gṛhakāryam)}]{\sk{गृहकार्यम्} (gṛhakāryam) — homework}

\label{lesson:SA-C255-grhakaryam}



\begin{quote}
Before the new one: say the Sanskrit for knowledge, then the Sanskrit for practice.
\end{quote}

\subsection*{You'll want to know: \sk{गृहकार्यम्}}
\textbf{\sk{गृहकार्यम्}} — \emph{gṛhakāryam} — \textquotedblleft{}homework\textquotedblright{}.

\begin{grammarlens}[title={What you've built}]
One more word for this chapter.
\end{grammarlens}

\subsection*{Your turn}
\begin{itemize}
\raggedright
  \item \emph{Say it:} \emph{gṛhakāryam}
  \item \emph{Say it:} \emph{gṛhakāryam}, once more
  \item \emph{Say it:} say \emph{gṛhakāryam}
  \item \emph{Recall:} say the Sanskrit for knowledge, then the Sanskrit for practice, then say \emph{gṛhakāryam} again
\end{itemize}

\begin{tcolorbox}[breakable,colback=teal!4,colframe=teal!35!black,title={Before you move on}]
What does \sk{गृहकार्यम्} mean? (\textquotedblleft{}Homework\textquotedblright{}.) Say it once more.
\end{tcolorbox}
\section[\texorpdfstring{\sk{उत्तरम्} (uttaram)}{उत्तरम् (uttaram)}]{\sk{उत्तरम्} (uttaram) — an answer}

\label{lesson:SA-C255-uttaram}



\begin{quote}
Before the new one: say the Sanskrit for practice, then the Sanskrit for homework.
\end{quote}

\subsection*{You'll want to know: \sk{उत्तरम्}}
\textbf{\sk{उत्तरम्}} — \emph{uttaram} — \textquotedblleft{}an answer\textquotedblright{}.

\begin{grammarlens}[title={What you've built}]
One more word for this chapter.
\end{grammarlens}

\subsection*{Your turn}
\begin{itemize}
\raggedright
  \item \emph{Say it:} \emph{uttaram}
  \item \emph{Say it:} \emph{uttaram}, once more
  \item \emph{Say it:} say \emph{uttaram}
  \item \emph{Recall:} say the Sanskrit for practice, then the Sanskrit for homework, then say \emph{uttaram} again
\end{itemize}

\begin{tcolorbox}[breakable,colback=teal!4,colframe=teal!35!black,title={Before you move on}]
What does \sk{उत्तरम्} mean? (\textquotedblleft{}An answer\textquotedblright{}.) Say it once more.
\end{tcolorbox}
